\documentclass[11pt]{article}
\usepackage[margin=1in]{geometry}
\usepackage{amsmath, amssymb, amsthm}
\usepackage{hyperref}

\newtheorem{theorem}{Theorem}
\newtheorem{proposition}[theorem]{Proposition}
\newtheorem{lemma}[theorem]{Lemma}
\newtheorem{corollary}[theorem]{Corollary}
\theoremstyle{remark}
\newtheorem{remark}[theorem]{Remark}

\DeclareMathOperator{\Tr}{Tr}
\DeclareMathOperator{\Hilb}{Hilb}
\DeclareMathOperator{\SYT}{SYT}
\newcommand{\Sym}{\mathrm{Sym}}
\newcommand{\ftilde}{\widetilde{f}}
\newcommand{\qe}{q_e}
\newcommand{\zetae}{\zeta_e}
\newcommand{\NN}{\mathbb{N}}
\newcommand{\ZZ}{\mathbb{Z}}
\newcommand{\QQ}{\mathbb{Q}}
\newcommand{\CC}{\mathbb{C}}

\title{Support of the Fake-Degree Schur Sum at a Root of Unity}
\author{Clio}
\date{August 6, 2026}

\begin{document}
\maketitle

\begin{abstract}
For $n \ge 1$ and $e \ge 2$, define
\[
  \qe := \sum_{\lambda \vdash n} \ftilde_\lambda(\zetae)\, s_\lambda \;\in\; \Lambda_n \otimes \ZZ[\zetae],
\]
where $\ftilde_\lambda(t)$ is the fake degree of $\lambda$ and $\zetae$ is a primitive $e$-th root of unity. We prove that $\qe$ is supported on power-sum monomials $p_\mu$ with
\[
  \mu = (e^{\lfloor n/e \rfloor},\, \nu), \quad \nu \vdash n \bmod e.
\]
The proof is a two-line consequence of the Chevalley--Molien formula for the graded character of the $S_n$-coinvariant algebra. As a corollary, $\qe$ factors as $p_e^{\lfloor n/e\rfloor} \cdot F_{n \bmod e}$ for a computable $F_{n \bmod e} \in \Lambda_{n \bmod e} \otimes \ZZ[\zetae]$.
\end{abstract}

\section{Statement}

Let $\Lambda_n$ denote the space of homogeneous symmetric functions of degree $n$ over $\QQ$. For $\lambda \vdash n$, the \emph{fake degree} of $\lambda$ is
\[
  \ftilde_\lambda(t) := \sum_{T \in \SYT(\lambda)} t^{\mathrm{maj}(T)} \;=\; t^{n(\lambda)}\,\frac{[n]_t!}{\prod_{c \in \lambda}[h(c)]_t},
\]
where $n(\lambda) = \sum_i (i-1)\lambda_i$ and $h(c)$ is the hook length of the cell $c$. Equivalently, $\ftilde_\lambda(t)$ is the graded multiplicity of the irreducible $V_\lambda$ in the coinvariant algebra of $S_n$.

Fix $e \ge 2$ and let $\zetae$ be a primitive $e$-th root of unity. Set
\[
  \qe := \sum_{\lambda \vdash n} \ftilde_\lambda(\zetae)\, s_\lambda \in \Lambda_n \otimes \ZZ[\zetae].
\]

\begin{theorem}\label{thm:support}
Write $\qe = \sum_\mu c_\mu\, p_\mu$ in the power-sum basis. Then $c_\mu \ne 0$ if and only if $\mu$ has the form
\[
  \mu = (e^{\lfloor n/e \rfloor},\, \nu), \qquad \nu \vdash n \bmod e.
\]
\end{theorem}

Equivalently, writing $\langle \cdot, \cdot\rangle$ for the Hall inner product, $\langle \qe, p_\mu\rangle = 0$ unless $\mu$ has $\lfloor n/e\rfloor$ parts equal to $e$ and the remaining parts sum to $n \bmod e$.

\section{Chevalley--Molien Preliminaries}

Let $V = \CC x_1 + \cdots + \CC x_n$ carry the natural $S_n$-action by permutation of coordinates, and let $S = \CC[V^*] = \CC[x_1, \ldots, x_n]$ with the induced graded action. The graded character
\[
  \Tr(g \mid S)(t) := \sum_{d \ge 0} \Tr(g \mid S_d)\, t^d
\]
is computed by Molien's formula:
\[
  \Tr(g \mid S)(t) = \frac{1}{\det_V(1 - t g)}.
\]
For $g \in S_n$ of cycle type $\mu = (\mu_1, \ldots, \mu_\ell)$, the permutation matrix $g$ has, on each cycle of length $\mu_i$, eigenvalues $\{\zeta_{\mu_i}^{0}, \ldots, \zeta_{\mu_i}^{\mu_i - 1}\}$. Since $\prod_{k=0}^{m-1}(1 - t\zeta_m^k) = 1 - t^m$, we obtain
\begin{equation}\label{eq:molien}
  \Tr(g \mid S)(t) = \frac{1}{\prod_{j \in \mu}(1 - t^j)}.
\end{equation}

Chevalley's theorem gives $S \cong S^{S_n} \otimes_\CC H$ as graded $S_n$-modules, where $H := S / (S^{S_n}_{+})$ is the coinvariant algebra and $S^{S_n} = \CC[e_1, \ldots, e_n]$ (with $\deg e_i = i$) acts trivially. Consequently
\[
  \Tr(g \mid S)(t) = \Tr(g \mid S^{S_n})(t) \cdot \Tr(g \mid H)(t) = \frac{1}{\prod_{i=1}^n (1 - t^i)} \cdot \Tr(g \mid H)(t),
\]
where we used that $g$ acts trivially on $S^{S_n}$ so $\Tr(g \mid S^{S_n})(t) = \Hilb(S^{S_n}; t) = \prod_{i=1}^n (1-t^i)^{-1}$. Solving for $\Tr(g \mid H)$ and combining with \eqref{eq:molien} gives
\begin{equation}\label{eq:cheva-molien}
  \Tr(g \mid H)(t) = \frac{\prod_{i=1}^n (1 - t^i)}{\prod_{j \in \mu}(1 - t^j)}.
\end{equation}

On the other hand, $H$ decomposes as $\bigoplus_{\lambda \vdash n} V_\lambda \otimes W_\lambda(t)$ with $W_\lambda(t)$ a graded multiplicity space of graded dimension $\ftilde_\lambda(t)$, so
\begin{equation}\label{eq:H-decomp}
  \Tr(g \mid H)(t) = \sum_{\lambda \vdash n} \ftilde_\lambda(t)\, \chi^\lambda(\mu).
\end{equation}

Combining \eqref{eq:cheva-molien} and \eqref{eq:H-decomp} yields the identity we need.

\begin{lemma}[Chevalley--Molien for fake degrees]\label{lem:cm}
For every $n \ge 1$ and $\mu \vdash n$,
\[
  \Psi_\mu(t) \;:=\; \sum_{\lambda \vdash n} \ftilde_\lambda(t)\, \chi^\lambda(\mu) \;=\; \frac{\prod_{i=1}^n (1 - t^i)}{\prod_{j \in \mu}(1 - t^j)}
\]
as elements of $\ZZ[t]$.
\end{lemma}

\begin{proof}
Immediate from \eqref{eq:cheva-molien} and \eqref{eq:H-decomp}. The ratio on the right is a polynomial (with integer coefficients) because it is the graded trace on the finite-dimensional coinvariant algebra $H$.
\end{proof}

\section{Proof of Theorem \ref{thm:support}}

For $\lambda \vdash n$, the Schur expansion in the power-sum basis is
\[
  s_\lambda = \sum_{\mu \vdash n} z_\mu^{-1} \chi^\lambda(\mu)\, p_\mu, \qquad z_\mu := \prod_i i^{m_i(\mu)}\, m_i(\mu)!,
\]
where $m_i(\mu)$ is the multiplicity of $i$ in $\mu$. Substituting into the definition of $\qe$ and swapping sums,
\[
  \qe = \sum_{\mu \vdash n} z_\mu^{-1} \biggl(\underbrace{\sum_{\lambda \vdash n} \ftilde_\lambda(\zetae)\, \chi^\lambda(\mu)}_{=\, \Psi_\mu(\zetae)\, \text{by Lemma \ref{lem:cm}}}\biggr) p_\mu.
\]
Hence
\begin{equation}\label{eq:cmu}
  c_\mu = z_\mu^{-1}\, \Psi_\mu(\zetae) = \frac{1}{z_\mu} \cdot \frac{\prod_{i=1}^n (1 - \zetae^i)}{\prod_{j \in \mu}(1 - \zetae^j)}.
\end{equation}

The factor $(1 - \zetae^m)$ vanishes if and only if $e \mid m$. Since $\Psi_\mu(t)$ is a polynomial (Lemma \ref{lem:cm}), we may factor $\Psi_\mu(t) = \Phi_e(t)^{v_e(\mu)}\, Q_\mu(t)$ with $Q_\mu(\zetae) \ne 0$, where
\begin{equation}\label{eq:vem}
  v_e(\mu) \;=\; \bigl|\{i \in [1,n] : e \mid i\}\bigr| - \bigl|\{j \in \mu : e \mid j\}\bigr| \;=\; \lfloor n/e \rfloor - k_e(\mu),
\end{equation}
with $k_e(\mu) := |\{j \in \mu : e \mid j\}|$.

\begin{lemma}\label{lem:vanish}
$c_\mu = 0$ if and only if $k_e(\mu) < \lfloor n/e \rfloor$.
\end{lemma}
\begin{proof}
From \eqref{eq:cmu}, $c_\mu = 0$ iff $\Psi_\mu(\zetae) = 0$ iff $v_e(\mu) > 0$, iff $k_e(\mu) < \lfloor n/e \rfloor$. We need $v_e(\mu) \ge 0$: parts of $\mu$ divisible by $e$ are each at least $e$, so their number $k_e(\mu)$ satisfies $e \cdot k_e(\mu) \le n$, hence $k_e(\mu) \le \lfloor n/e \rfloor$.
\end{proof}

\begin{proof}[Proof of Theorem \ref{thm:support}]
By Lemma \ref{lem:vanish}, $c_\mu \ne 0$ iff $k_e(\mu) = \lfloor n/e \rfloor$. Let $k := \lfloor n/e \rfloor$. Suppose $k_e(\mu) = k$; write the parts of $\mu$ divisible by $e$ as $e a_1, \ldots, e a_k$ with $a_i \ge 1$. Their sum is $e \sum_i a_i$, and the remaining parts sum to $n - e\sum_i a_i \ge 0$, so $\sum_i a_i \le n/e$, i.e., $\sum_i a_i \le k$. Combined with $\sum_i a_i \ge k$, we get $\sum_i a_i = k$, so every $a_i = 1$.

Thus the parts of $\mu$ divisible by $e$ are exactly $\lfloor n/e\rfloor$ copies of $e$. Let $\nu$ be the partition formed by the remaining parts. Then $|\nu| = n - ek = n \bmod e < e$, so all parts of $\nu$ are strictly less than $e$ (in particular not multiples of $e$). Hence $\mu = (e^{\lfloor n/e\rfloor}, \nu)$ with $\nu \vdash n \bmod e$.

Conversely, if $\mu = (e^{\lfloor n/e\rfloor}, \nu)$ with $\nu \vdash n \bmod e$, then $k_e(\mu) = \lfloor n/e\rfloor$ (the parts of $\nu$ are $\le n \bmod e < e$, hence not multiples of $e$), so $v_e(\mu) = 0$ and $c_\mu \ne 0$.
\end{proof}

\section{Corollary: Product Structure and Explicit Formula}

Since $p_{(e^k, \nu)} = p_e^k \cdot p_\nu$ in $\Lambda$, Theorem \ref{thm:support} immediately gives:

\begin{corollary}\label{cor:factor}
With $k := \lfloor n/e \rfloor$ and $\rho := n \bmod e$,
\[
  \qe = p_e^k \cdot F_e,
\]
where $F_e \in \Lambda_\rho \otimes \ZZ[\zetae]$ has power-sum expansion
\[
  F_e = \sum_{\nu \vdash \rho} c_{(e^k, \nu)}\, p_\nu.
\]
\end{corollary}

We can compute the coefficients $c_{(e^k, \nu)}$ in closed form. For $\mu = (e^k, \nu)$,
\[
  z_\mu = e^k \cdot k! \cdot z_\nu, \qquad \Psi_\mu(\zetae) = \frac{\prod_{i=1}^n (1 - \zetae^i)}{(1 - \zetae^e)^k \prod_{j \in \nu}(1 - \zetae^j)}.
\]
Both numerator and denominator vanish to order $k$ at $\zetae$; factoring out $\Phi_e(t)^k$ and evaluating,
\[
  \lim_{t \to \zetae} \frac{\prod_{m=1}^k(1 - t^{em})}{(1 - t^e)^k} = \prod_{m=1}^k \Bigl(\lim_{t \to \zetae}\frac{1 - t^{em}}{1 - t^e}\Bigr) = \prod_{m=1}^k m = k!,
\]
using $\lim_{u \to 1}(1-u^m)/(1-u) = m$ with $u = t^e$. Writing $\alpha_r := 1 - \zetae^r$ for $r \in \{1, \ldots, e-1\}$ and setting $b_r := k + \mathbf{1}[r \le \rho]$, the surviving factors give
\begin{equation}\label{eq:closed}
  c_{(e^k, \nu)} = \frac{1}{e^k\, z_\nu}\, \prod_{r=1}^{e-1} \alpha_r^{\,b_r - m_r(\nu)}.
\end{equation}

\begin{remark}
Formula \eqref{eq:closed} matches the values computed in Section \ref{sec:verif}:
\begin{itemize}
\item At $(n, e) = (5, 3)$: $k = 1$, $\rho = 2$, so $b_1 = b_2 = 2$. For $\nu = (2)$: $m_1 = 0, m_2 = 1$, so $c_{(3,2)} = \tfrac{1}{3 \cdot 2}\, \alpha_1^2 \alpha_2 = \tfrac{\alpha_1^2 \alpha_2}{6}$. For $\nu = (1,1)$: $m_1 = 2, m_2 = 0$, so $c_{(3,1,1)} = \tfrac{1}{3 \cdot 2}\, \alpha_2^2 = \tfrac{\alpha_2^2}{6}$.
\item At $(n, e) = (20, 9)$: $k = 2$, $\rho = 2$, and $\Psi_{(9,9,2)}(\zeta_9)$ evaluates to $162 - 162\zeta_9$, matching the $2! \cdot 81 = 162$ prefactor.
\end{itemize}
\end{remark}

\section{Computational Verification}\label{sec:verif}

Two independent scripts verify Theorem \ref{thm:support}.

\paragraph{(a) Chevalley--Molien polynomial identity.} For every $n \le 7$ and every $\mu \vdash n$, we verified $\sum_\lambda \ftilde_\lambda(t) \chi^\lambda(\mu) = \prod(1-t^i)/\prod(1-t^j)$ as polynomials in $\ZZ[t]$. This checks Lemma \ref{lem:cm} at $\sum_{n \le 7} p(n) = 30$ pairs $(n, \mu)$.

\paragraph{(b) Support of $\qe$.} For each of the pairs
\[
  (n, e) \in \{(5, 3), (6, 3), (7, 5), (8, 3), (9, 3), (10, 4), (11, 3), (15, 4)\},
\]
we computed the coefficient $c_\mu$ of $p_\mu$ in $\qe$ over all $\mu \vdash n$ (total $\sum p(n) = 7 + 11 + 15 + 22 + 30 + 42 + 56 + 176 = 359$ tests), reducing modulo $\Phi_e(t)$. Every $\mu$ not of the shape $(e^{\lfloor n/e\rfloor}, \nu)$ gave $c_\mu = 0$, and every $\mu$ of that shape gave $c_\mu \ne 0$. Sample:
\begin{center}
\begin{tabular}{rlrl}
\hline
$(n, e)$ & $\mu$ & & $\Psi_\mu(\zetae)$ (writing $t := \zetae$) \\
\hline
$(5, 3)$ & $(3, 2)$ & & $3 - 3t$ \\
$(5, 3)$ & $(3, 1, 1)$ & & $3 + 3t$ \\
$(7, 5)$ & $(5, 2)$ & & $5 - 5t$ \\
$(7, 5)$ & $(5, 1, 1)$ & & $5 + 5t$ \\
$(8, 3)$ & $(3, 3, 2)$ & & $18 - 18t$ \\
$(8, 3)$ & $(3, 3, 1, 1)$ & & $18 + 18t$ \\
$(10, 4)$ & $(4, 4, 2)$ & & $32 - 32t$ \\
$(10, 4)$ & $(4, 4, 1, 1)$ & & $32 + 32t$ \\
\hline
\end{tabular}
\end{center}
The prefactors ($3, 5, 18, 32$) match $k!\, \prod \alpha_r^{b_r - m_r(\nu)}$ from \eqref{eq:closed}, reduced modulo $\Phi_e$.

\paragraph{(c) The (20, 9) instance.} Yesterday's probe (2026-08-05) verified over all 627 partitions of $20$: the only $\mu$ with $c_\mu \ne 0$ are $(9, 9, 2)$ with $\Psi_\mu(\zeta_9) = 162 - 162\, \zeta_9$ and $(9, 9, 1, 1)$ with $\Psi_\mu(\zeta_9) = 162 + 162\, \zeta_9$. Both match \eqref{eq:closed} with $k = 2$, $2! \cdot 81 = 162$.

\section{Remarks and Cross-References}

\begin{remark}[Springer-regular case as a special case]
When $e$ divides $n$ and $e \in \{$degrees of $S_n\} = \{1, 2, \ldots, n\}$ (i.e., $e \le n$), the setup is the Kra\'skiewicz--Weyman regularity case: $\lfloor n/e\rfloor = n/e$, $n \bmod e = 0$, so the only $\nu$ is empty and $\qe = c \cdot p_e^{n/e}$ for a single coefficient $c$. This recovers Springer's regular-element formula.
\end{remark}

\begin{remark}[Connection to cyclic sieving]
The polynomial $\Psi_\mu(t)$ is exactly the graded character of the $S_n$-coinvariant ring at a permutation of cycle type $\mu$. Reiner, Stanton and White \cite{RSW04} use exactly the ratio $\prod(1-t^i)/\prod(1-t^j)$ (evaluated at roots of unity) as the sieving polynomial in a foundational example of the cyclic sieving phenomenon. Theorem \ref{thm:support} is a straightforward consequence in that language: $\Psi_\mu(\zetae)$ counts fixed points of $\langle \sigma\rangle$ (with $\sigma$ a fixed element of order $e$ in $S_n$) acting on the coset space $S_n/C(\mu)$, and this count is zero unless $\sigma$ can be conjugated inside a maximal torus normalizer for cycle-type $\mu$, i.e., unless $\mu$ contains $\lfloor n/e\rfloor$ copies of $e$.
\end{remark}

\begin{remark}[The Albion Verschiebung route is not needed]
The PROVE plan (2026-08-06) suggested closing the conjecture via Albion's Theorem 1.3 on Verschiebung acting on Schur functions \cite{Alb26}: apply $\varphi_e$ to both sides of a putative product identity $\qe = c\, p_e^k F_{n\bmod e}$, using $\varphi_e s_\lambda = 0$ unless the $e$-core of $\lambda$ is empty. This route works but is heavier: it requires identifying the $e$-quotient structure that makes the sum collapse, and it recovers only the same order-of-vanishing statement that falls out of the Chevalley--Molien formula in one line. The Chevalley--Molien path is more direct because it goes through the coinvariant algebra, where the Verschiebung structure is already \emph{built in} to the isomorphism $H \cong \CC[S_n]/(\ldots)$ via the graded character formula.
\end{remark}

\begin{remark}[Application to composite-$d$ divisibility]
Theorem \ref{thm:support} unlocks the mechanism behind Theorem B of \cite{Cli26} (unconditional $\Phi_9 \mid m_\pi$ at $(r, k) = (5, 4)$): for any triple $(r, k, e)$ where the factorisation-count obstruction fires (no partition of $r$ can produce $\lfloor n/e\rfloor$ copies of $e$ in $p_\lambda[h_k]$), we obtain $\Phi_e \mid m_\pi(t)$ for all $\pi \vdash r$. This anchors the composite-$d$ paper alongside Theorems A + B.
\end{remark}

\begin{thebibliography}{9}
\bibitem{RSW04} V.~Reiner, D.~Stanton, and D.~White,
\emph{The cyclic sieving phenomenon},
J.~Combin.~Theory Ser.~A \textbf{108} (2004), 17--50.

\bibitem{Alb26} K.~Albion,
\emph{Character-factorisations and Verschiebung operators on symmetric functions},
arXiv:2501.18520 (Jan 2025).

\bibitem{Cli26} Clio,
\emph{Composite-$d$ divisibility for wreath byproduct polynomials: $k=2$ closure and unconditional $\Phi_9$ at $(5,4)$},
\texttt{proofs/2026-08-05-composite-d-k2-and-negative.tex}.

\bibitem{Hum90} J.~E.~Humphreys,
\emph{Reflection Groups and Coxeter Groups},
Cambridge Univ.~Press, 1990. (\S 3.6 for the Chevalley/Steinberg formula on coinvariants.)

\bibitem{KW01} W.~Kra\'skiewicz and J.~Weyman,
\emph{Algebra of coinvariants for a reflection group},
J.~Algebra \textbf{240} (2001), 745--762.
\end{thebibliography}

\end{document}
